\documentclass{article}
\input{../../lib.sty}

\title{\texttt{fn score\_candidates}}
\author{\MichaelShoemateAuthor}
\begin{document}
\maketitle


\contrib

Proves soundness of \rustdoc{transformations/quantile\_score\_candidates/fn}{score\_candidates}
in \asOfCommit{mod.rs}{f5bb719}.
\texttt{score\_candidates} returns a score for each candidate passed in,
where the score is the distance between the candidate and the ideal alpha-quantile.

\section{Hoare Triple}
\subsection*{Precondition}
\subsubsection*{Compiler-verified}
\begin{itemize}
    \item Generic \texttt{TIA} (input atom type) is a type with trait \texttt{PartialOrd}.
\end{itemize}

\subsubsection*{Caller-verified}
\begin{itemize}
    \item \texttt{x} has fewer than $2^{64}$ elements.
    \item The elements of \texttt{x} and \texttt{candidates} belong to a common totally ordered set.
    \item \texttt{candidates} is strictly increasing
    \item $\texttt{alpha\_num} \le \texttt{alpha\_den}$
    \item $\texttt{size\_limit} \cdot \texttt{alpha\_den} < 2^{64}$
\end{itemize}


\subsection*{Pseudocode}
\label{sec:python-pseudocode}
\lstinputlisting[language=Python,firstline=2,escapechar=`]{./pseudocode/score_candidates.py}


\subsection*{Postcondition}

\begin{theorem}
    Let $C$ denote \texttt{candidates} and let $l$ denote \texttt{size\_limit}.
    For each index $i$ in $\{1, \ldots, |C|\}$,

    \begin{equation*}
        \mathrm{score\_candidates}_i(x, C, \alpha_{num}, \alpha_{den}, l) = |(\alpha_{den}-\alpha_{num}) \cdot \min(\#(x < C_i), l) - \alpha_{num} \cdot \min(\#(x > C_i), l)|
    \end{equation*}

    where $\#(x < C_i)$ is the number of elements in $x$ less than $C_i$,
    $\#(x > C_i)$ is the number of elements in $x$ greater than $C_i$.
\end{theorem}

\begin{proof}
    The function breaks down into two parts:
    \begin{itemize}
        \item Compute histograms, where the edges are the candidates.
        \item Uses the histograms to compute scores.
    \end{itemize}

    The histograms are initialized at zero, with one more bin than candidates,
    since the bins start at $-\infty$ and end at $+\infty$.

    The bins in \texttt{hist\_ro} are closed on the left, and open on the right.
    The bins in \texttt{hist\_lo} are open on the left, and closed on the right.

    On line \ref{hist-lo-increment}, \texttt{idx\_lt} is the number of candidates strictly less than \texttt{x\_i}, identifying its left-open bin.
    On line \ref{hist-ro-increment}, \texttt{idx\_eq} is the number of candidates less than or equal to \texttt{x\_i}, identifying its right-open bin.
    Since candidates are strictly increasing, the equality search adds one if a candidate equals \texttt{x\_i}, and zero otherwise.

    Line \ref{n-sum} computes the total number of elements in the data as each element in \texttt{x} incremented a bin in \texttt{hist\_lo}.
    No arithmetic thus far can overflow due to the precondition that \texttt{x} has fewer than $2^{64}$ elements.

    Notice that there is one more bin than candidates.
    With \texttt{n}, the last bin in each of the histograms (the number of elements beyond the largest candidate),
    is not needed to compute the score.
    Therefore, the last bin is discarded on lines \ref{pop-ro} and \ref{pop-lo}.

    The scores can now be computed in one linear pass over the histograms on line \ref{zip-hist}.
    The number of elements less than the candidate \texttt{lt} is the cumulative sum of \texttt{hist\_ro},
    and the number of elements less than or equal to the candidate \texttt{le} is the cumulative sum of \texttt{hist\_lo}.

    The number of elements greater than the candidate \texttt{gt} is then simply \texttt{n - le} on line \ref{gt}.

    To ensure that the score computation does not overflow, the counts are bounded by \texttt{size\_limit} on line \ref{lt-gt-lim}.

    By the precondition that $\texttt{alpha\_num} \le \texttt{alpha\_den}$,
    then both \texttt{alpha\_den - alpha\_num} and \texttt{alpha\_num} are less than or equal to \texttt{alpha\_den}.
    Using this, and the precondition that $\texttt{size\_limit} \cdot \texttt{alpha\_den} < 2^{64}$,
    then the computation of the score on line \ref{score} is guaranteed to not overflow.

    The computation on line \ref{score} directly satisfies the postcondition.
\end{proof}


\end{document}
